\section{Пример сгенерированного конспекта}

Ниже приведён фрагмент конспекта, сформированного RAG-сервисом для одного из анонимных
профилей ученика (идентификатор \texttt{archive\_u5}) на основе выборки решений с портала
ЕГЭПрорыв. Исходные данные о попытках решения задач получены из деперсонализированной
выгрузки (36 учеников; удалены ФИО, адреса электронной почты и иные идентификаторы).
Полный текст конспекта в формате Markdown сохранён в репозитории разработки как часть
архива экспериментальных прогонов.

\subsection*{Введение (фрагмент)}

Давайте разберём два важных задания первой части ЕГЭ. Задание №10 --- прикладная математика
(подстановка в формулу), а №12 --- исследование функции с помощью производной. В обоих
случаях ответом должно быть целое число или конечная десятичная дробь.

\subsection*{Что важно запомнить (фрагмент)}

\begin{enumerate}
    \item \textbf{Производная логарифма: $(\ln x)' = \frac{1}{x}$.}
    Логарифм растёт всё медленнее по мере увеличения $x$. Если в №12 видишь функцию с $\ln$,
    помни про ОДЗ: аргумент логарифма всегда $> 0$.
    \item \textbf{Производная частного:}
    $\left(\frac{f}{g}\right)' = \frac{f'g - fg'}{g^2}$.
    В формуле производной частного в числителе всегда стоит минус перед слагаемым $fg'$.
    \item \textbf{Правило <<догона>> (преследование с форой):}
    $t = \frac{S_{\mathrm{фора}}}{v_2 - v_1}$ --- время догона при разнице скоростей.
\end{enumerate}

\subsection*{Типичные ошибки (фрагмент)}

\begin{enumerate}
    \item \textbf{Игнорирование ограничений (ОДЗ) в текстовых задачах.}
    При решении квадратного уравнения в №10 могут получиться два корня; скорость, время и длина
    не могут быть отрицательными --- выбирай физически осмысленный корень.
    \item \textbf{Забытый минус в производной частного.}
    Путать знак с производной произведения нельзя: в числителе всегда $f'g - fg'$.
\end{enumerate}

\subsection*{Разбор примера: задание №12 (фрагмент)}

\textbf{Условие.} Найдите наименьшее значение функции
$f(x) = x^2 - 8x + 19$ на отрезке $[1; 5]$.

\textbf{Решение.}
\begin{enumerate}
    \item Производная: $f'(x) = 2x - 8$.
    \item Критическая точка: $2x - 8 = 0 \Rightarrow x = 4$.
    \item Проверка: $4 \in [1; 5]$.
    \item Значения: $f(1) = 12$, $f(5) = 4$, $f(4) = 3$.
    \item Наименьшее из $\{12; 4; 3\}$ равно $3$.
\end{enumerate}

\textbf{Ответ:} $3$.

\subsection*{Чеклист перед ответом (фрагмент)}

\begin{itemize}
    \item В ответе нет букв, корней, $\pi$ и обыкновенных дробей (если этого требует формат задания).
    \item В задаче №12 вычислены значения функции на концах отрезка и в критических точках внутри отрезка.
    \item Единицы измерения в текстовой задаче №10 приведены к одной системе.
\end{itemize}
